\documentclass[11pt]{article}
\usepackage[margin=1in]{geometry}
\usepackage{amsmath,amssymb,amsthm,mathtools}
\usepackage{booktabs}
\usepackage{array}
\usepackage{fontspec}
\usepackage{fancyvrb}
\usepackage{needspace}
\newfontfamily\leanfont{DejaVuSansMono.ttf}
\newfontfamily\symfont{STIXTwoMath-Regular.otf}
\newcommand{\sym}[1]{{\symfont #1}}
\usepackage[hidelinks]{hyperref}

\newtheorem{theorem}{Theorem}
\newtheorem{proposition}[theorem]{Proposition}
\newtheorem{corollary}[theorem]{Corollary}
\theoremstyle{definition}
\newtheorem{definition}{Definition}
\theoremstyle{remark}
\newtheorem{remark}{Remark}

\newcommand{\lean}[1]{\texttt{\small #1}}
\newcommand{\R}{\mathbb{R}}
\newcommand{\sech}{\operatorname{sech}}
\newcommand{\diag}{\operatorname{diag}}

\title{Saturated Signed Ripple Dynamics on Directed Networks:\\
Machine-Checked Theory and a Preregistered Empirical Test}
\author{Ruhai Huang\\ \small Wanyouzhiliu Tech}
\date{Preprint v0.1.1 --- November 2026}

\setlength{\emergencystretch}{2em}
\begin{document}
\maketitle

\begin{abstract}
We study a minimal model of signed influence (``ripples'') on a directed network,
$r(t+1)=D(t)\,r(t)+W(t)\tanh r(t)+s(t)$, with a signed, unbounded state, retention $D$, a signed and
\emph{unnormalised} influence matrix $W$, and saturating transmission. We assemble known results from
recurrent-network stability, single-neuron bifurcation theory and spectral graph theory into one signed
regime map, and machine-check all of it in Lean~4 with Mathlib. The results are: global boundedness
under arbitrary switching; a common-weight contraction certificate, equivalent for constant
coefficients to $\rho(|D|+|W|)<1$; a scalar regime map indexed by the signed gain $K=d+w$, including the
critical law $t\,x(t)^2\to3/(2w)$ at $K=1$; the inequality $\rho(W)\le\rho(|W|)$, with equality under
structural balance; exact identifiability criteria; and a global formula for the input--response gain.
The development has 50 theorems and 63 lemmas (1{,}907 lines, no \texttt{sorry}, standard axioms only),
and the formalisation refuted three of our own informal claims. Two lightweight-preregistered fits of a
reduced-form version on the Bitcoin OTC signed network found no detectable propagation gain over
self-dynamics, with a minimum detectable effect of about $5\%$. Distrust edges did reverse sign, as
balance theory predicts. We specify the data and experiments needed for a decisive test.
\end{abstract}

\section{Introduction}

Many systems carry \emph{signed} influence along \emph{directed} links. Trust and distrust ratings, approval
and disapproval among agents, excitatory and inhibitory couplings, and supportive versus adversarial
messages in multi-agent software systems are all examples. A common modelling instinct is to propagate a
scalar ``ripple'' linearly, $r(t+1)=Ar(t)+s(t)$, and to read off stability from the spectral radius of
$A$. This has three well-known problems. Linear growth is unbounded, while real influence saturates.
Per-time-step spectral conditions do not control switched systems. And row-normalising $A$ to restore
stability erases the very gain one wishes to measure.

This paper develops the smallest model we found that avoids these problems while remaining
falsifiable, and tests it. Our contributions are the following.

\begin{enumerate}
\item \textbf{Model (Section~\ref{sec:model}).} A saturated, signed, directed ripple equation in which the
      local signed gain $K=d+w$ is a dynamical parameter that may be negative or exceed one, rather
      than a score clipped to $[0,1]$.
\item \textbf{Theory (Section~\ref{sec:theory}).} Global boundedness under switching, a common-weight
      contraction certificate and its spectral equivalent, a complete scalar regime map indexed by $K$
      including the exact critical rate, the formally verified half of a frustration-attenuation result
      linking structural balance to dynamics, identifiability criteria, and a global gain formula.
\item \textbf{Machine-checked proofs (Section~\ref{sec:formal}).} Every theorem stated as ``verified'' is
      proved in Lean~4/Mathlib. Two proofs avoid machinery that Mathlib lacks. The spectral certificate
      avoids the Perron--Frobenius theorem, and the gain formula is proved from a second-order Taylor
      bound instead of the implicit function theorem.
\item \textbf{Self-correction.} The formalisation refuted three claims of our own earlier draft. We list
      them in Table~\ref{tab:corrections}, because we regard documented refutations as part of the
      evidence.
\item \textbf{Preregistered empirical test (Section~\ref{sec:empirical}).} Two rounds on a public signed
      network. The results are mostly negative for the propagation term, and we report them as such.
\end{enumerate}

\paragraph{Relation to prior work and novelty.} The model is a discrete-time additive recurrent network
of Hopfield/echo-state type~\cite{hopfield1984,jaeger2001}. Sufficient conditions for global stability of
such networks via (diagonally) scaled norms and similarity transformations are
classical~\cite{jin1994,buehner2006,yildiz2012}. Contraction
analysis more broadly is due to~\cite{lohmiller1998}. The inequality $\rho(W)\le\rho(|W|)$ and its
equality case are Wielandt's theorem~\cite{wielandt1950,hornjohnson}. Structural balance goes back
to~\cite{heider1946,harary1953,cartwright1956}, and dynamics on signed networks with antagonistic
interactions were studied by~\cite{altafini2013}. Walk-based spectral measures of balance, such as
$\operatorname{tr}e^{A}/\operatorname{tr}e^{|A|}$~\cite{estrada2014}, and their
critiques~\cite{kirkley2019} are close relatives of our ratio $\eta$. Bifurcations of a single
self-coupled model neuron, including period doubling and hysteresis, were analysed
by~\cite{pasemann1993}; our scalar regime map is an instance of that analysis. Lean~4 formalisations of
Hopfield networks, Boltzmann machines and the Perron--Frobenius theorem for irreducible nonnegative
matrices were given by~\cite{cipollina2025}. Critical $t^{-1/2}$ relaxation for maps
$x\mapsto x-ax^3+O(x^5)$ is classical asymptotics~\cite{debruijn1958}. We do \emph{not} claim these
individual mathematical facts as new. Our contribution is to assemble them into a single, signed regime
map for an influence model. We also supply complete kernel-checked proofs, some in a more general form
than the textbook statements (time-varying coefficients, non-diagonal $D$), document where formalisation corrected informal claims, and
confront the model with preregistered data.

\section{Model}\label{sec:model}

\begin{definition}[Saturated signed ripple equation]
For $n$ nodes, let
\begin{equation}\label{eq:model}
  r(t+1)=D(t)\,r(t)+W(t)\,\phi\big(r(t)\big)+s(t),
\end{equation}
where $\phi$ acts componentwise. In the base model $\phi=\tanh$, $D=\diag(d_i)$ with $|d_i|<1$,
$W_{ij}\in\R$ is the signed influence of $j$ on $i$, and $s$ is a bounded external input. The state
$r_i$ is signed and unbounded. Positive and negative values denote the two directions of a
channel-specific reference axis.
\end{definition}

The retention term is decomposed as $d_i=\beta_i-\lambda_i\kappa_i$: an intrinsic retention $\beta_i$ minus
a dissipation $\lambda_i$ that scales with an externally controllable precision $\kappa_i$. Several
channels per node are handled by flattening $r$ into $\R^{nc}$ with block matrices. All results below hold
in that dimension, and Kronecker-separable couplings are a special case.

\paragraph{Why saturation.} Since $\tanh r=r-r^3/3+O(r^5)$, we have $W\tanh r\approx Wr-Wr^3/3$. Truncating
at the cubic term instead diverges in discrete time at moderate amplitude: for $x\mapsto0.9x-0.5x^3$
started at $x_0=2.5$, the fourth iterate has magnitude $9.4\times10^{15}$. The saturated form keeps the small-amplitude behaviour and removes the divergence.

\paragraph{The signed gain $K$.} Linearising at the origin gives the Jacobian $J=D+W$. For a single node,
$f(x)=dx+w\tanh x$ and $K:=f'(0)=d+w$. $K$ is not normalised. Its sign and size select the regime
(Theorem~\ref{thm:regimes}).

\section{Theory}\label{sec:theory}

Throughout, $\|\cdot\|$ is the $\ell^\infty$ norm, $|A|$ is the entrywise absolute value, and
$\rho$ is the spectral radius. Lean identifiers of the verified statements are given in brackets.

\subsection{Boundedness and contraction}

\begin{theorem}[Global boundedness under arbitrary switching; \lean{P1\_global\_bound}]\label{thm:bound}
Let $D(t),W(t)$ be arbitrary real matrices with row sums of absolute values $\le d<1$ and $\le w$
respectively. Let $\|s(t)\|\le S$ and $|\phi|\le 1$. Then every solution of~\eqref{eq:model} satisfies
\[
  \|r(t)\|\le d^t\|r(0)\|+(w+S)\frac{1-d^t}{1-d}.
\]
\end{theorem}

The proof is an elementary one-line induction. The point is that the bound does not involve the
spectrum of $W$. Influence can be amplified, up to saturation, but
cannot explode. Note that $D$ need not be diagonal.

\begin{definition}[Common weight certificate]
$(v,c)$ with $v>0$ and $0\le c<1$ is a common weight for $(D(t),W(t))_t$ if
$\sum_j\big(|D_{ij}(t)|+|W_{ij}(t)|\big)v_j\le c\,v_i$ for all $t,i$.
\end{definition}

\begin{theorem}[Global contraction; \lean{P2\_contraction}, \lean{P2\_exists\_equilibrium},
\lean{P2\_unique\_equilibrium}, \lean{P2\_global\_attraction}]\label{thm:contraction}
Let $\phi$ be $1$-Lipschitz and suppose a common weight $(v,c)$ exists. Then any two solutions satisfy
$|r_i(t)-q_i(t)|\le C_0\,c^t v_i$ whenever $|r_i(0)-q_i(0)|\le C_0 v_i$. If $D,W,s$ are constant, there
is a unique equilibrium, and it is globally attracting.
\end{theorem}

\begin{theorem}[Spectral form; \lean{weight\_iff\_rho\_lt\_one}, \lean{P2\_tanh\_of\_rho}]\label{thm:spectral}
For a nonnegative $M\in\R^{n\times n}$ with $n\ge1$, a common weight exists, i.e.\ $Mv\le cv$ with $v>0$ and
$c<1$, if and only if $\rho(M)<1$. Hence $\rho(|D|+|W|)<1$ suffices for Theorem~\ref{thm:contraction}
in the time-invariant case.
\end{theorem}

\begin{remark}
For switched systems, the relevant object is a \emph{common} certificate, not the per-step spectrum.
Section~\ref{sec:switch} gives a machine-checked counterexample. Our proof of the ``if'' direction
avoids Perron--Frobenius. It uses Gelfand's formula to find $N$ with $\|M^N\|<r^N$ for
$\rho(M)<r<1$, and sets $v=\sum_{k<N}(M/r)^k\mathbf 1$.
\end{remark}

\subsection{The scalar regime map}

Let $f(x)=dx+w\tanh x$ with $|d|<1$ and $K=d+w$. A uniform remainder $|f(x)-Kx|\le|w|\,|x|^3$ underlies
the local statements.

\needspace{14\baselineskip}
\begin{theorem}[Regimes indexed by $K$]\label{thm:regimes}
\leavevmode
\begin{center}
\small\upshape
\begin{tabular}{@{}l p{0.47\linewidth} p{0.31\linewidth}@{}}
\toprule
$K$ & behaviour & Lean \\
\midrule
$K<-1$ & stable period-2 orbit $\pm p$ (non-decaying oscillation) & \lean{P3\_period\_two} \\
$-1<K<0$ & locally, alternating-sign geometric decay & \lean{P3\_alternating\_decay} \\
$K=0$ & linear response vanishes; $|f(x)|\le|w||x|^3$ & \lean{P3\_decoupled} \\
$0<K<1$ & locally, sign-preserving geometric decay & \lean{P3\_monotone\_decay} \\
$K=1$ & global convergence to $0$, slower than any geometric rate & \lean{P3\_K1\_converges}, \lean{P3\_K1\_slow} \\
$K>1$ & origin unstable; stable symmetric fixed points $\pm p$ & \lean{P3\_origin\_unstable}, \lean{P3\_bistable} \\
\bottomrule
\end{tabular}
\end{center}
\end{theorem}

The bifurcations at $K=\pm1$ (pitchfork and period doubling) are those of the classical single-neuron
map~\cite{pasemann1993}. What the formal development adds is complete proofs, including the existence
and stability of the nontrivial fixed points and of the $2$-cycle.

\begin{theorem}[Exact critical slowing; \lean{P3\_K1\_rate}, \lean{P3\_K1\_rate\_pos}]\label{thm:critical}
Let $K=1$ and $|d|<1$, so $w=1-d>0$. If the orbit never hits $0$ exactly, then
\[
  t\,x(t)^2\longrightarrow\frac{3}{2w},\qquad\text{equivalently}\qquad \sqrt t\,|x(t)|\to\sqrt{3/(2w)}.
\]
For $0\le d<1$ and $x_0\neq0$ the orbit never hits $0$, so the conclusion holds unconditionally.
\end{theorem}

\begin{proof}[Proof idea]
Two derivative-sign arguments give the sandwich $x^3/3-2x^5/15\le x-\tanh x\le x^3/3$ for $x\ge0$.
Hence $1/x(t+1)^2-1/x(t)^2\to 2w/3$, and the Cesàro mean gives the result.
\end{proof}

Theorem~\ref{thm:critical} makes a prediction that is not visible from $K$ alone. Two nodes at the same
critical gain $K=1$ relax at different speeds, and recovery is slower when more of the gain is carried
by linear retention $d$ than by saturating influence $w$. Numerically, with $2\times10^6$ steps from
$x_0=1$, $t\,x^2$ equals $3.000005$, $1.499997$ and $15.0004$ for $(d,w)=(0.5,0.5),(0,1),(0.9,0.1)$. The
predicted values are $3$, $1.5$ and $15$.

\begin{proposition}[Sign-determined basins and their failure; \lean{P3\_basin}, \lean{P3\_basin\_fails\_neg\_d}]
If $K>1$ and $0\le d<1$, every $x_0>0$ converges to $+p$ and every $x_0<0$ converges to $-p$. For $d<0$ this
fails: with $d=-0.9$ and $w=2$ (so $K=1.1$) we have $f(10)<0$, so a large positive perturbation can
end in the negative state.
\end{proposition}

\subsection{Frustration attenuates propagation}

\begin{theorem}[\lean{P4\_rho\_le\_rho\_abs}, \lean{P4\_balanced\_eq}, \lean{P4\_eta\_le\_one}, \lean{P4\_eta\_balanced}]\label{thm:frustration}
For every real matrix $W$, $\rho(W)\le\rho(|W|)$. If $W$ is structurally balanced or antibalanced,
i.e.\ $W_{ij}=\theta\,\sigma_i|W_{ij}|\sigma_j$ with $\sigma_i,\theta\in\{\pm1\}$, then
$\rho(W)=\rho(|W|)$. Consequently the frustration attenuation factor
$\eta:=\rho(W)/\rho(|W|)$ satisfies $\eta\le1$, with $\eta=1$ under balance.
\end{theorem}

No irreducibility is needed for the inequality. The proof uses $|W^k|\le|W|^k$ entrywise and Gelfand's
formula. The converse, that equality for irreducible $W$ forces balance up to a diagonal unitary
similarity, is Wielandt's theorem~\cite{wielandt1950}. It is \emph{not} formalised here, because it requires
Perron--Frobenius theory, which Mathlib~4.34 does not provide. An external Lean~4 formalisation of
Perron--Frobenius for irreducible nonnegative matrices~\cite{cipollina2025} offers a route to close this
gap. Theorem~\ref{thm:frustration} is therefore the formally verified \emph{half} of the interface between
the topology of a signed graph and the dynamics. The substantive claim, that frustrated cycles
strictly reduce linear propagation gain, rests on the classical theorem and is not machine-checked here.

\paragraph{Frustration count is not the right statistic.}
We drew 249 random strongly connected
weighted digraphs with 3--7 nodes. A balanced signing $S|W|S$ of each gave $\eta=1$ in all 249 cases.
Random signings that were frustrated and not antibalanced (234 cases) gave $\eta<1$ in 232, with median
$\eta=0.68$. The two exceptions were periodic graphs with complex-phase balance, consistent with
Wielandt's equality case. However, the Spearman correlation between $\eta$ and the frustration index
(the minimum number of sign flips needed for balance, divided by the number of edges) was only
$-0.105$ ($p=0.11$). The spectral ratio $\eta$, not a count of ``bad'' edges, should be used.
$\eta$ belongs to the family of walk-based balance indices~\cite{estrada2014}. As a ratio of spectral radii
it isolates the dominant long-walk contribution, which is the quantity that controls linear
propagation gain. Whether it shares the large-network bias reported for such
indices~\cite{kirkley2019} is open.

\subsection{Identifiability}

The model is linear in its parameters. For node $i$, $r_i(t+1)$ is regressed on
$[r_i(t),\tanh r(t),s_i(t)]$ with design matrix $Z$.

\begin{theorem}[\lean{P5\_identifiable\_iff}, \lean{P5\_ols\_exact}, \lean{P5\_ols\_error},
\lean{P5\_diag\_separable}, \lean{P5\_diag\_weak}, \lean{P5\_diag\_gram}, \lean{P5\_kappa\_identify}]\label{thm:ident}
\leavevmode
\begin{enumerate}
\item Parameters are unique iff $\det(Z^\top Z)\ne0$. Least squares then recovers them exactly, and with
      noise $\varepsilon$ the error is exactly $(Z^\top Z)^{-1}Z^\top\varepsilon$.
\item Retention $d$ and self-influence $W_{ii}$ are separable from two samples $x_1,x_2$ iff
      $x_1\ne0$, $x_2\ne0$ and $|x_1|\ne|x_2|$. The key lemma is that $\tanh x/x$ is strictly decreasing on $(0,\infty)$.
\item Along the direction $(d+c,\,W_{ii}-c)$ predictions change by at most $|c||x|^3$. Correspondingly, the
      quadratic form of the Gram matrix of $[x,\tanh x]$ along $(1,-1)$ is at most $\sum_t x_t^6$. At small
      amplitude, separating $d$ from $W_{ii}$ is a sixth-order effect, and noise is amplified along
      that direction.
\item In $d=\beta-\lambda\kappa$, $\beta$ and $\lambda$ are identifiable iff $\kappa$ takes at least two values.
\end{enumerate}
\end{theorem}

Item~3 says that saturation is what makes self-dynamics identifiable. Item~4 is a hard experimental
design requirement: testing ``higher precision implies stronger dissipation'' requires externally
varied $\kappa$.

\subsection{Input--response gain}

\begin{theorem}[Global gain formula; \lean{P6\_invertible}, \lean{P6\_equilibrium\_map}, \lean{P6\_lipschitz}, \lean{P6\_gain}]\label{thm:gain}
Assume a common weight $(v,c)$ for constant $D,W$, with $\phi=\tanh$. Then:
\begin{enumerate}
\item $I-D-W\diag(\sech^2x)$ is invertible at every $x$.
\item The equilibrium map $s\mapsto r^*(s)$ is well defined and Lipschitz, with constant
      \[ L=\Big(\sum_j v_j^{-1}\Big)\Big(\sum_j v_j\Big)\Big/(1-c). \]
\item $r^*$ is Fréchet differentiable at every $s$, with
      \[ G(s)=\frac{\partial r^*}{\partial s}=\big(I-D-W\diag(\sech^2 r^*(s))\big)^{-1}. \]
\end{enumerate}
\end{theorem}

$G_{ij}$ is the signed amount of ripple ultimately produced at $i$ by a unit excitation at $j$. This is
the model's behavioural quantity. The proof uses the second-order remainder
$|\tanh(a+\delta)-\tanh a-\sech^2a\,\delta|\le2\delta^2$ rather than the implicit function theorem. The
formula holds at every $s$ because the common-weight condition makes the equilibrium globally unique.

\begin{proposition}[Cooperative networks; \lean{P6\_cooperative}]
If in addition $D,W\ge0$ entrywise, then $0\le G\le(I-D-W)^{-1}$ entrywise: saturation can only reduce gains.
\end{proposition}

In signed networks this monotonicity fails. With $d=0$, $w=-0.9$ and constant input $s=1.5$, saturation
raises the gain from the linear value $0.526$ to $0.686$. Across 100 random systems the analytic $G$ matched finite differences to
$2.4\times10^{-9}$, and the median ratio of saturated to linear total gain was $0.70$.

\subsection{Switching: a machine-checked counterexample}\label{sec:switch}

\begin{proposition}[Switching counterexample]
Let $A_1=\begin{psmallmatrix}0.3&3\\0&0.3\end{psmallmatrix}$ and
$A_2=\begin{psmallmatrix}0.3&0\\3&0.3\end{psmallmatrix}$. Each has the single eigenvalue $0.3$. Alternating
$r(t+1)=A_{t\bmod 2}\,r(t)$ from $(1,1)$ gives $r_2(2k)\ge9^k$, so the trajectory is unbounded. The
common-weight certificate correctly refuses this family: any common weight needs $c\ge3.3$. The
saturated version $r(t+1)=0.3\,r+W(t)\tanh r$, with $W(t)$ switching arbitrarily between the two
off-diagonal parts, obeys $\|r(t)\|\le0.3^t\|r(0)\|+3(1-0.3^t)/0.7$.
(Lean: \lean{SW\_single\_eigen}, \lean{SW\_linear\_diverges}, \lean{SW\_no\_common\_weight},
\lean{SW\_saturated\_bounded}.)
\end{proposition}

\section{Formalisation}\label{sec:formal}

All results above marked with Lean identifiers are proved in a single file, \texttt{KRippleV2.lean},
using Lean~4~\cite{moura2021} and Mathlib~\cite{mathlib2020}.

\begin{center}
\begin{tabular}{@{}ll@{}}
\toprule
Toolchain & Lean \texttt{v4.34.1}, Mathlib (8{,}529 modules) \\
Size & 1{,}907 lines; 50 theorems, 63 lemmas \\
Audit & no \texttt{sorry}; \texttt{\#print axioms} lists only \texttt{propext}, \texttt{Classical.choice}, \texttt{Quot.sound} \\
Build & \texttt{lake build}, about 40\,s with a cached Mathlib, zero warnings \\
SHA-256 & \texttt{\footnotesize b9f8fa17691531c0625648de6530cfd6859c3fdf108b0a8c0a038879cbe37ce3} \\
\bottomrule
\end{tabular}
\end{center}

The formal statements are often more general than the informal draft. Theorem~\ref{thm:bound} allows
arbitrary $D(t)$ and any $|\phi|\le1$, and Theorem~\ref{thm:contraction} requires only a $1$-Lipschitz
$\phi$. Mathlib did not contain the Lipschitz property or the derivative of $\tanh$, so both were proved
along the way. Appendix~\ref{app:lean} reproduces the key definitions and theorem statements verbatim,
extracted from the source by script, so that readers can check the informal statements against the
formal ones.

\begin{table}[h]
\centering
\caption{Claims of our own earlier informal draft that formalisation refuted or sharpened.}\label{tab:corrections}
\begin{tabular}{@{}p{0.42\linewidth}p{0.52\linewidth}@{}}
\toprule
Informal claim & Outcome \\
\midrule
For $K>1$ the final state is determined by the sign of the initial perturbation. & False for $d<0$
(explicit counterexample). True for $d\ge0$, and for small perturbations. \\ \addlinespace
Saturation reduces the input--response gain. & Not a theorem in signed networks
(counterexample above). True entrywise in cooperative networks. \\ \addlinespace
$d$ and $W_{ii}$ separate only above $|r|\approx0.1$. & No mathematical threshold exists: any two nonzero
samples of different magnitude separate them. The threshold is empirical, set by noise against an
$O(|x|^3)$ signal, as quantified in Theorem~\ref{thm:ident}(3). \\
\bottomrule
\end{tabular}
\end{table}

\paragraph{Not formalised.}
\begin{itemize}
\item The strict direction of Wielandt's theorem, which needs Perron--Frobenius.
\item Complex regimes (rotation and invariant circles for $\rho>1$, a Neimark--Sacker
      bifurcation~\cite{kuznetsov}).
\item Statistical consistency of the estimators under dependent, noisy data.
\end{itemize}
Each is a substantial independent project.

\section{Synthetic numerical checks}\label{sec:numerics}

A script with seed 20261118 (numpy~2.5.2, scipy~1.18.1, about 15\,s) cross-checks the theory on random systems
with up to 12 nodes. Theorem~\ref{thm:bound} was violated 0 times in 300 random switched systems with
initial magnitudes of order $10^4$, and the observed maximum reached $99.6\%$ of the bound, so the bound
is tight. Uniqueness of equilibria under the spectral condition failed 0 times in 300 systems. The
$K$-regime prediction matched the observed behaviour in $222/222$ cells of a $(d,w)$ grid, including
32 cells where only local guarantees apply. In the identifiability check, with 4 nodes, noise $0.01$ and
$T=400$, the errors of $d$ and $W_{ii}$ were about $65$ at input amplitude $0.01$, where the condition
number was $8.6\times10^3$. At amplitude $1$ they fell to $10^{-3}$. Off-diagonal entries had errors below $0.01$ at every amplitude. With constant $\kappa$ the design had rank $3/4$, giving
$\hat\beta=0.31$ and $\hat\lambda=-0.24$ against true values $0.9$ and $0.5$. With varying $\kappa$ the
rank was $4/4$, giving $\hat\beta=0.897$ and $\hat\lambda=0.499$.

\section{Preregistered empirical test: Bitcoin OTC}\label{sec:empirical}

\paragraph{What is tested.} We do not fit equation~\eqref{eq:model} edge by edge. Each round fits a
\emph{reduced form}: a node's own state plus four global coefficients on signed neighbour aggregates
(incoming and outgoing, trust and distrust). The test therefore asks whether signed neighbour
aggregates carry predictive information, which is a necessary but not sufficient consequence of the
model. Fitted values of $K$ are reduced-form self-gains, not structural estimates.

\paragraph{Data.} The Bitcoin OTC signed trust network~\cite{kumar2016,snap} contains 35{,}592 timestamped
ratings in $\{-10,\dots,+10\}$. We binned them by week and split them in time: the first $70\%$ of weeks
for training and the remainder for testing.

\paragraph{Protocol.} Before each round of fitting, an analysis plan was written and its SHA-256 hash
recorded (\texttt{PLAN.sha256}, \texttt{PLAN\_V2.sha256}). This is \emph{lightweight} preregistration,
because the same analyst wrote the plan and ran the fit~\cite{nosek2018}. The hashes were recorded in
our own repository and were not externally timestamped. They document intent but cannot by themselves
prove that the plans preceded the analyses. Uncertainty was estimated with
a whole-week block bootstrap (1{,}000 resamples).

\paragraph{Round 1: direct observation model.} The state $r_i(t)$ was the sum of ratings received in
week $t$, divided by 10. We compared the saturated signed model against constant, persistence, linear,
unsigned and permuted-network baselines. Four of five preregistered hypotheses failed. On all 470\,k
node-weeks, $99.6\%$ of which receive no rating the following week, the saturated model had a test MSE
$5.6\%$ higher than persistence. Diagnosis showed that the failure lay in the observation model, which
conflated ``no event'' with ``zero ripple''. The hypothesis that held was sign reversal through distrust
edges: the coefficient was $+0.026$ for incoming trust and $-0.069$ for incoming distrust, essentially
unchanged after controlling for node degree.

\paragraph{Round 2: two-part observation model.} An event model (logistic regression on activity) was
combined with a magnitude model fitted on event rows, using a decaying-accumulation state with a time
constant fixed in advance. Because this design was informed by round~1 and reuses its test period, the
evidence is \emph{weaker} than a fresh test. A confirmatory run on a second, untouched dataset (Bitcoin
Alpha) is specified in the plan but has not yet been run.

\begin{table}[h]
\centering
\caption{Round 2 results. MSE is computed on the 1{,}761 test rows in which an event occurred, and AUC
is for predicting a net-negative rating given an event.}
\begin{tabular}{@{}lccc@{}}
\toprule
Model & event-row MSE & combined MSE & conditional AUC \\
\midrule
B1 self-dynamics only & 0.18761 & 0.000799 & \textbf{0.790} \\
B2 linear network & 0.18626 & 0.000800 & 0.761 \\
B3 saturated signed (ours) & 0.18679 & 0.000806 & 0.761 \\
B4 unsigned & 0.18763 & 0.000804 & 0.738 \\
B5 permuted network & 0.18631 & 0.000803 & 0.761 \\
\bottomrule
\end{tabular}
\end{table}

\paragraph{Findings.}
\begin{itemize}
\item The two-part model fixed the observation problem. Combined MSE was $3.4\%$ below round-1
      persistence (95\% CI excludes 0), and the event AUC was $0.95$.
\item The network propagation term gave \emph{no} detectable gain over self-dynamics: $-0.4\%$, with a
      CI containing 0. Saturation, sign and true topology were likewise not significantly better than
      their controls.
\item This negative result has limited power. The 95\% CI of the event-row MSE difference,
      $[-0.0098,+0.0084]$, spans about $\pm5\%$ of the event-row MSE ($\approx0.187$). Propagation
      effects smaller than about $5\%$ therefore cannot be excluded.
\item Distrust sign reversal held again, with coefficients $+0.021$ and $-0.106$.
\item Sign early-warning was significant (AUC $0.76$, CI $0.73$--$0.80$), but it came from the node's own
      history, since baseline B1 reached $0.79$.
\item The fitted self-gain was $K=d+w_0=0.54$, which lies in the dissipative regime.
\end{itemize}

\paragraph{Interpretation.} In this sparse, low-frequency rating network, what survives is saturating
self-dissipation, with $K\approx0.5$, together with an event-intensity model. Cross-node ripples are
too weak to yield predictive gain, although their coefficients carry the sign pattern predicted by
balance. This does not refute the propagation part of the theory, but it does not support it either.
Testing it requires denser, faster interaction data and an analysis powered to detect effects of
about $1\%$ (Section~\ref{sec:agenda}).

\section{Limitations}

The model is a hypothesis about influence dynamics. The theorems are about the model, not about any
real system. The observation map from events to $r$ is the main open modelling problem, as round~1
showed. The empirical evidence is a single dataset, a single analyst, lightweight preregistration and a
reused test period in round~2. The estimates use global coefficients without node heterogeneity. The
choice of $\tanh$ is one of many: every bounded, $1$-Lipschitz, odd saturation preserves the main
results, and the functional form should be chosen from data.

\section{Open problems and collaboration agenda}\label{sec:agenda}

We are seeking academic or laboratory partners for the following work. Each item has a falsification
criterion that is fixed in advance.

\begin{enumerate}
\item[\textbf{E1}] \textbf{Confirmatory replication.} Preregister publicly (e.g.\ on OSF), then run the
      frozen round-2 specification unchanged on
      the Bitcoin Alpha network and on further signed temporal networks. Failure criterion: event-row
      MSE of B3 not below B1, with the 95\% CI excluding 0.
\item[\textbf{E2}] \textbf{Dense interaction data for the propagation term.} Use timestamped, signed or
      stance-labelled interactions with fast propagation, such as multi-agent LLM message traces,
      reply and forward cascades with stance labels, or signed hyperlink networks between online
      communities. Hypothesis: network terms improve held-out prediction over self-dynamics, and
      $\eta$ predicts attenuation better than frustrated-edge counts.
\item[\textbf{E3}] \textbf{Interventional design for $d=\beta-\lambda\kappa$.} Run experiments in which a
      precision or attention variable $\kappa$ is externally varied. This is required by
      Theorem~\ref{thm:ident}(4) and cannot be replaced by observational data.
\item[\textbf{E4}] \textbf{Critical-slowing signature.} Test whether nodes estimated near $K\approx1$ show
      $t^{-1/2}$ relaxation after shocks, with a constant ordered by $w$ as Theorem~\ref{thm:critical}
      predicts, in contrast to generic early-warning indicators~\cite{scheffer2009}.
\item[\textbf{F1}] \textbf{Formal mathematics.} Formalise Perron--Frobenius and the strict Wielandt
      direction in Lean/Mathlib, and establish statistical consistency under mixing.
\end{enumerate}

\paragraph{What a partner would provide.}
\begin{itemize}
\item Access to suitable datasets (E2, E3) under the partner's data governance.
\item Computing resources for node-heterogeneous fitting on $10^6$--$10^8$ node-time rows, large
      bootstrap ensembles, multi-agent simulation for E2, and CPU time for Mathlib-scale formal
      developments (F1).
\item Independent preregistration and analysis, which would upgrade the evidence from single-analyst
      to independent replication.
\end{itemize}

\paragraph{What we provide.}
\begin{itemize}
\item The Lean library.
\item Frozen analysis plans with hashes and scripted, seed-fixed pipelines.
\item A commitment to report negative results.
\end{itemize}
The theory requires little computation by itself. Resources are needed for data scale and independent
verification.

\section*{Data, code and AI-use statement}
The Lean development, all numerical and analysis scripts, their outputs, and the preregistered plans
with their hashes are public at \url{https://github.com/853948826-bit/k-ripple-v2} (code under the MIT
licence, this manuscript under CC BY 4.0) and archived on Zenodo, \href{https://doi.org/10.5281/zenodo.24568466}{doi:10.5281/zenodo.24568466}.
The Bitcoin OTC data are public~\cite{snap} and are not redistributed; the scripts verify the file's
SHA-256 before use. The formal proofs, analysis scripts and the drafting of this
manuscript were carried out with the assistance of an AI coding agent (Claude Code, Anthropic). Every
theorem labelled as verified is checked by the Lean kernel, and every number reported is produced by a
script in the accompanying repository. Responsibility for the content rests with the author.

\appendix
\section{Key Lean statements}\label{app:lean}
The following definitions (with their bodies) and theorem statements (up to \texttt{:=}) are extracted
verbatim from \texttt{KRippleV2.lean} (SHA-256 above) by the script \texttt{gen\_lean\_appendix.py}. Proofs
are omitted. \texttt{rho} is the spectral radius of the complexified matrix, valued in $[0,\infty]$. The
constant $C_0$ of Theorem~\ref{thm:contraction} is named \texttt{K} in the Lean statement.

% generated by gen_lean_appendix.py from KRippleV2.lean sha256 b9f8fa17691531c0625648de6530cfd6859c3fdf108b0a8c0a038879cbe37ce3
\begin{Verbatim}[fontsize=\footnotesize,formatcom=\leanfont,commandchars=¦«»]
def RowSumLe (A : Matrix (Fin n) (Fin n) ℝ) (a : ℝ) : Prop :=
  ∀ i, ∑ j, |A i j| ≤ a

def CommonWeight (D W : ℕ → Matrix (Fin n) (Fin n) ℝ) (v : Fin n → ℝ) (c : ℝ) : Prop :=
  ∀ t i, ∑ j, (|D t i j| + |W t i j|) * v j ≤ c * v i

def IsTraj (D W : ℕ → Matrix (Fin n) (Fin n) ℝ) (s : ℕ → Fin n → ℝ) (φ : ℝ → ℝ)
    (r : ℕ → Fin n → ℝ) : Prop :=
  ∀ t, r (t + 1) = (D t).mulVec (r t) + (W t).mulVec (fun i => φ (r t i)) + s t

def stepMap (D W : Matrix (Fin n) (Fin n) ℝ) (s : Fin n → ℝ) (φ : ℝ → ℝ)
    (x : Fin n → ℝ) : Fin n → ℝ :=
  D.mulVec x + W.mulVec (fun i => φ (x i)) + s

noncomputable def sech2 (x : ℝ) : ℝ := 1 / Real.cosh x ^ 2

noncomputable def rho (M : Matrix (Fin n) (Fin n) ℝ) : ℝ≥0∞ :=
  spectralRadius ℂ (M.map (algebraMap ℝ ℂ))

def absMat (W : Matrix (Fin n) (Fin n) ℝ) : Matrix (Fin n) (Fin n) ℝ :=
  Matrix.of fun i j => |W i j|

noncomputable def scalarMap (d w : ℝ) (x : ℝ) : ℝ := d * x + w * Real.tanh x

noncomputable def linJ (D W : Matrix (Fin n) (Fin n) ℝ) (x : Fin n → ℝ) :
    Matrix (Fin n) (Fin n) ℝ :=
  D + W * Matrix.diagonal (fun j => sech2 (x j))

theorem P1_global_bound (D W : ℕ → Matrix (Fin n) (Fin n) ℝ) (s : ℕ → Fin n → ℝ)
    (φ : ℝ → ℝ) (r : ℕ → Fin n → ℝ) {d w S : ℝ}
    (hd0 : 0 ≤ d) (hd1 : d < 1) (hw : 0 ≤ w)
    (hD : ∀ t, RowSumLe (D t) d) (hW : ∀ t, RowSumLe (W t) w) (hs : ∀ t, ‖s t‖ ≤ S)
    (hφ : ∀ y, |φ y| ≤ 1)
    (hr : ∀ t, r (t + 1) = (D t).mulVec (r t) + (W t).mulVec (fun i => φ (r t i)) + s t)
    (t : ℕ) :
    ‖r t‖ ≤ d ^ t * ‖r 0‖ + (w + S) * (1 - d ^ t) / (1 - d)

theorem P2_contraction {D W : ℕ → Matrix (Fin n) (Fin n) ℝ} {s : ℕ → Fin n → ℝ}
    {φ : ℝ → ℝ} {r q : ℕ → Fin n → ℝ} {v : Fin n → ℝ} {c K : ℝ}
    (hφ : ∀ a b, |φ a - φ b| ≤ |a - b|) (hr : IsTraj D W s φ r) (hq : IsTraj D W s φ q)
    (hc0 : 0 ≤ c) (hM : CommonWeight D W v c) (hK0 : 0 ≤ K)
    (hK : ∀ i, |r 0 i - q 0 i| ≤ K * v i) :
    ∀ t i, |r t i - q t i| ≤ K * c ^ t * v i

theorem weight_iff_rho_lt_one [NeZero n] {M : Matrix (Fin n) (Fin n) ℝ}
    (hM : ∀ i j, 0 ≤ M i j) :
    (∃ v : Fin n → ℝ, (∀ i, 0 < v i) ∧
      ∃ c : ℝ, 0 ≤ c ∧ c < 1 ∧ ∀ i, ∑ j, M i j * v j ≤ c * v i) ↔ rho M < 1

theorem P3_K1_rate {d w : ℝ} (hd : |d| < 1) (hK : d + w = 1) (x₀ : ℝ)
    (hne : ∀ t, (scalarMap d w)^[t] x₀ ≠ 0) :
    Tendsto (fun t : ℕ => (t : ℝ) * ((scalarMap d w)^[t] x₀) ^ 2) atTop (¦sym«𝓝» (3 / (2 * w)))

theorem P4_rho_le_rho_abs (W : Matrix (Fin n) (Fin n) ℝ) : rho W ≤ rho (absMat W)

theorem P4_balanced_eq {W : Matrix (Fin n) (Fin n) ℝ} {σ : Fin n → ℝ} {θ : ℝ}
    (hσ : ∀ i, σ i ^ 2 = 1) (hθ : θ ^ 2 = 1) (hW : ∀ i j, W i j = θ * σ i * |W i j| * σ j) :
    rho W = rho (absMat W)

theorem P5_diag_separable (x₁ x₂ : ℝ) :
    x₁ * Real.tanh x₂ - x₂ * Real.tanh x₁ ≠ 0 ↔ x₁ ≠ 0 ∧ x₂ ≠ 0 ∧ |x₁| ≠ |x₂|

theorem P6_gain (hv : ∀ i, 0 < v i) (hc1 : c < 1)
    (hM : CommonWeight (fun _ => D) (fun _ => W) v c)
    (rs : (Fin n → ℝ) → Fin n → ℝ) (hrs : ∀ s, rs s = stepMap D W s Real.tanh (rs s))
    (s : Fin n → ℝ) :
    HasFDerivAt rs (Matrix.toLin' (1 - linJ D W (rs s))⁻¹).toContinuousLinearMap s

theorem P6_cooperative (hv : ∀ i, 0 < v i) (hc1 : c < 1)
    (hM : CommonWeight (fun _ => D) (fun _ => W) v c)
    (hD : ∀ i j, 0 ≤ D i j) (hW : ∀ i j, 0 ≤ W i j) (x : Fin n → ℝ) :
    ∀ i j, 0 ≤ (1 - linJ D W x)⁻¹ i j ∧ (1 - linJ D W x)⁻¹ i j ≤ (1 - (D + W))⁻¹ i j
\end{Verbatim}


\begin{thebibliography}{99}\small
\bibitem{altafini2013} C.~Altafini. Consensus problems on networks with antagonistic interactions.
  \emph{IEEE Trans.\ Automatic Control}, 58(4):935--946, 2013.
\bibitem{buehner2006} M.~Buehner and P.~Young. A tighter bound for the echo state property.
  \emph{IEEE Trans.\ Neural Networks}, 17, 2006.
\bibitem{cartwright1956} D.~Cartwright and F.~Harary. Structural balance: a generalization of Heider's theory.
  \emph{Psychological Review}, 63(5):277--293, 1956.
\bibitem{cipollina2025} M.~Cipollina, M.~Karatarakis and F.~Wiedijk. Formalized Hopfield networks and
  Boltzmann machines. arXiv:2512.07766, 2025; LPAR-26, 2026.
\bibitem{debruijn1958} N.~G.~de~Bruijn. \emph{Asymptotic Methods in Analysis}. North-Holland, 1958.
\bibitem{estrada2014} E.~Estrada and M.~Benzi. Walk-based measure of balance in signed networks:
  detecting lack of balance in social networks. \emph{Physical Review E}, 90(4):042802, 2014.
\bibitem{harary1953} F.~Harary. On the notion of balance of a signed graph.
  \emph{Michigan Mathematical Journal}, 2(2):143--146, 1953.
\bibitem{heider1946} F.~Heider. Attitudes and cognitive organization.
  \emph{Journal of Psychology}, 21:107--112, 1946.
\bibitem{hopfield1984} J.~J.~Hopfield. Neurons with graded response have collective computational
  properties like those of two-state neurons. \emph{PNAS}, 81(10):3088--3092, 1984.
\bibitem{hornjohnson} R.~A.~Horn and C.~R.~Johnson. \emph{Matrix Analysis}, 2nd ed. Cambridge University Press, 2013.
\bibitem{jaeger2001} H.~Jaeger. The ``echo state'' approach to analysing and training recurrent neural
  networks. GMD Report 148, German National Research Center for Information Technology, 2001.
\bibitem{jin1994} L.~Jin, P.~N.~Nikiforuk and M.~M.~Gupta. Absolute stability conditions for discrete-time
  recurrent neural networks. \emph{IEEE Trans.\ Neural Networks}, 5(6):954--964, 1994.
\bibitem{kirkley2019} M.~Kirkley, G.~T.~Cantwell and M.~E.~J.~Newman. Balance in signed networks.
  \emph{Physical Review E}, 99:012320, 2019.
\bibitem{kumar2016} S.~Kumar, F.~Spezzano, V.~S.~Subrahmanian and C.~Faloutsos. Edge weight prediction in
  weighted signed networks. In \emph{Proc.\ IEEE ICDM}, pp.~221--230, 2016.
\bibitem{kuznetsov} Yu.~A.~Kuznetsov. \emph{Elements of Applied Bifurcation Theory}, 3rd ed.
  Applied Mathematical Sciences 112. Springer, 2004.
\bibitem{lohmiller1998} W.~Lohmiller and J.-J.~E.~Slotine. On contraction analysis for non-linear systems.
  \emph{Automatica}, 34(6):683--696, 1998.
\bibitem{mathlib2020} The mathlib Community. The Lean mathematical library. In \emph{Proc.\ CPP}, pp.~367--381, 2020.
\bibitem{moura2021} L.~de~Moura and S.~Ullrich. The Lean~4 theorem prover and programming language.
  In \emph{CADE-28}, LNCS 12699, pp.~625--635, 2021.
\bibitem{nosek2018} B.~A.~Nosek, C.~R.~Ebersole, A.~C.~DeHaven and D.~T.~Mellor. The preregistration revolution.
  \emph{PNAS}, 115(11):2600--2606, 2018.
\bibitem{pasemann1993} F.~Pasemann. Dynamics of a single model neuron. \emph{Int.\ J.\ Bifurcation and Chaos},
  3(2), 1993. doi:10.1142/S0218127493000210.
\bibitem{scheffer2009} M.~Scheffer et al. Early-warning signals for critical transitions.
  \emph{Nature}, 461:53--59, 2009.
\bibitem{snap} J.~Leskovec and A.~Krevl. SNAP Datasets: Stanford large network dataset collection.
  \url{http://snap.stanford.edu/data}, 2014.
\bibitem{wielandt1950} H.~Wielandt. Unzerlegbare, nicht negative Matrizen.
  \emph{Mathematische Zeitschrift}, 52:642--648, 1950.
\bibitem{yildiz2012} I.~B.~Yildiz, H.~Jaeger and S.~J.~Kiebel. Re-visiting the echo state property.
  \emph{Neural Networks}, 35:1--9, 2012.
\end{thebibliography}

\end{document}
